\chapter{Analyse et Conception du Système}
\label{ch:conception}

\section*{Introduction}
\phantomsection
\addcontentsline{toc}{section}{Introduction}

Ce chapitre décrit les choix de conception
du système. Le Chapitre~\ref{ch:etat_art}
a montré qu'à notre connaissance, aucun travail
publié ne combine détection temps réel,
anonymisation Privacy-by-Design et décision
adaptative sur un seul dispositif embarqué
dans le contexte marocain. Ce chapitre explique
comment on a structuré le système pour répondre
à ces trois contraintes simultanément.

L'architecture retenue s'articule autour de
six modules interconnectés. Chaque choix —
du modèle de détection au tracker, du format
d'export au noyau gaussien — est motivé par
des contraintes concrètes : le CPU ARM du
Raspberry Pi 5, la Loi 09-08, et l'exigence
de fonctionnement en temps réel.

\section{Analyse des besoins}

\subsection{Besoins fonctionnels}

Le système doit remplir cinq fonctions
qui s'enchaînent séquentiellement :

\begin{table}[H]
\centering
\caption{Besoins fonctionnels du système}
\begin{tabularx}{\textwidth}{clX}
\toprule
\textbf{BF} &
\textbf{Fonction} &
\textbf{Description} \\
\midrule
BF1 & Acquisition &
Lire un flux vidéo d'intersection depuis
un fichier local, frame par frame,
compatible avec le pipeline de traitement. \\

BF2 & Détection &
Identifier et localiser en temps réel
les objets de six classes : voitures,
bus, camions, motos, piétons et véhicules
d'urgence, avec un score de confiance
par détection. \\

BF3 & Anonymisation &
Flouter automatiquement chaque piéton
détecté, immédiatement après la détection
et avant tout traitement ultérieur,
conformément à la Loi 09-08. \\

BF4 & Décision &
Calculer la phase de feux à activer
selon les densités détectées par zone,
avec priorité pondérée aux piétons et
priorité absolue aux véhicules d'urgence. \\

BF5 & Visualisation &
Afficher le flux annoté, les statistiques
par zone et l'état des feux via un
dashboard accessible en réseau local. \\
\bottomrule
\end{tabularx}
\label{tab:besoins_fonctionnels}
\end{table}

\subsection{Besoins non-fonctionnels}

Au-delà des fonctions métier, plusieurs
contraintes qualitatives encadrent
l'architecture :

\begin{table}[H]
\centering
\caption{Besoins non-fonctionnels du système}
\begin{tabularx}{\textwidth}{clXl}
\toprule
\textbf{BNF} &
\textbf{Critère} &
\textbf{Exigence} &
\textbf{Cible} \\
\midrule
BNF1 & Temps réel &
Cadence suffisante pour une décision
cohérente avec le trafic réel &
$\geq$ 5 FPS \\

BNF2 & Précision &
Détection suffisante pour couvrir
la quasi-totalité des piétons &
mAP50 $\geq$ 75\% \\

BNF3 & Conformité légale &
Respect de la Loi 09-08 &
100\% \\

BNF4 & Embarquabilité &
Fonctionne sur CPU ARM sans GPU &
RAM $\leq$ 4 GB \\

BNF5 & Maintenabilité &
Modules indépendants, modifiables
séparément &
Architecture modulaire \\
\bottomrule
\end{tabularx}
\label{tab:besoins_non_fonctionnels}
\end{table}

\subsection{Contraintes matérielles}

Le Raspberry Pi 5 8GB est le dispositif
cible. Ses caractéristiques imposent
des choix architecturaux directs :

\begin{figure}[H]
\centering
\resizebox{0.92\textwidth}{!}{%
\begin{tikzpicture}[
  constraint/.style={
    draw=AccentOrange,
    fill=AccentOrange!10,
    rounded corners=4pt,
    align=center, font=\scriptsize,
    minimum width=4.2cm,
    minimum height=0.75cm},
  solution/.style={
    draw=AccentGreen,
    fill=AccentGreen!10,
    rounded corners=4pt,
    align=center, font=\scriptsize,
    minimum width=4.2cm,
    minimum height=0.75cm},
  arrow/.style={
    PrimaryBlue,->,thick}
]

\node[font=\small\bfseries, color=AccentOrange]
  at (0,0) {Contrainte matérielle};
\node[font=\small\bfseries, color=AccentGreen]
  at (6.5,0) {Solution retenue};

\node[constraint] (c1) at (0,-0.95)
  {CPU ARM Cortex-A76, pas de GPU};
\node[constraint] (c2) at (0,-1.85)
  {RAM partagée 8 GB};
\node[constraint] (c3) at (0,-2.75)
  {FP16 non supporté nativement};
\node[constraint] (c4) at (0,-3.65)
  {Bande passante réseau limitée};
\node[constraint] (c5) at (0,-4.55)
  {Risque de throttling thermique};

\node[solution] (s1) at (6.5,-0.95)
  {YOLO26n nano + ONNX Runtime};
\node[solution] (s2) at (6.5,-1.85)
  {Modèle nano FP32, ~9.8 MB};
\node[solution] (s3) at (6.5,-2.75)
  {Export ONNX FP32 retenu};
\node[solution] (s4) at (6.5,-3.65)
  {Streaming MJPEG compressé};
\node[solution] (s5) at (6.5,-4.55)
  {Pipeline séquentiel, pas de
  parallélisme forcé};

\foreach \c/\s in
  {c1/s1,c2/s2,c3/s3,c4/s4,c5/s5}{
  \draw[arrow] (\c.east) -- (\s.west);
}
\end{tikzpicture}}
\caption{Contraintes du Raspberry Pi 5
         et solutions architecturales}
\label{fig:constraints}
\end{figure}

\section{Architecture globale du système}

\subsection{Vue d'ensemble}

Le pipeline traite les frames de façon
séquentielle, de l'acquisition vidéo
jusqu'à la décision de phase de feux.
Chaque module reçoit les sorties du
précédent et transmet les siennes au
suivant.

\begin{figure}[H]
\centering
\resizebox{0.96\textwidth}{!}{%
\begin{tikzpicture}[
  font=\small,
  module/.style={
    draw=PrimaryBlue, fill=PrimaryBlue!8,
    rounded corners=6pt,
    minimum width=3.35cm,
    minimum height=1.05cm,
    align=center, very thick
  },
  privacy/.style={
    draw=AccentOrange, fill=AccentOrange!10,
    rounded corners=6pt,
    minimum width=3.35cm,
    minimum height=1.05cm,
    align=center, very thick
  },
  output/.style={
    draw=AccentGreen, fill=AccentGreen!8,
    rounded corners=6pt,
    minimum width=3.35cm,
    minimum height=1.05cm,
    align=center, very thick
  },
  arrow/.style={
    -{Latex[length=3mm]},
    line width=1.1pt,
    color=DarkGray,
    shorten >=1pt,
    shorten <=1pt},
  note/.style={
    font=\scriptsize, color=DarkGray,
    align=center, fill=white,
    inner sep=1.5pt}
]

\node[module] (m1) at (0,0)
  {M1\\Acquisition};
\node[module] (m2) at (5.0,0)
  {M2\\Détection YOLO};
\node[privacy] (m4) at (10.0,0)
  {M4\\Anonymisation};
\node[module] (m3) at (10.0,-2.05)
  {M3\\Tracking IoU};
\node[module] (m5) at (5.0,-2.05)
  {M5\\Décision MDP};
\node[output] (m6) at (0,-2.05)
  {M6\\Dashboard};

\draw[arrow] (m1) --
  node[note, above] {frames} (m2);
\draw[arrow] (m2) --
  node[note, above] {bboxes} (m4);
\draw[arrow] (m4) --
  node[note, right=3pt] {frame floutée}
  (m3);
\draw[arrow] (m3) --
  node[note, above=3pt, text width=1.4cm]
  {zones\\+ densités} (m5);
\draw[arrow] (m5) --
  node[note, above=3pt, text width=1.4cm]
  {phase\\+ durée} (m6);

\node[
  draw=AccentOrange,
  fill=AccentOrange!5,
  rounded corners=5pt,
  align=center,
  font=\scriptsize,
  text width=8.6cm
] at (5.0,-3.25)
{\textbf{Point clé :} M4 floute les personnes
avant le tracking et avant l'affichage.
Aucun pixel identifiable ne circule
au-delà de M4.};
\end{tikzpicture}}
\caption{Pipeline du système — M1 à M6
         \textit{(Source : auteur)}}
\label{fig:pipeline_system}
\end{figure}

\subsection{Source vidéo retenue}

Le pipeline utilise une vidéo du Bellevue
Traffic Video Dataset stockée localement
sur le Raspberry Pi 5. Ce choix permet
de rejouer exactement les mêmes séquences
pour comparer les modèles, les tailles
d'entrée et les paramètres de décision.
La calibration des polygones est faite
une fois sur une frame de référence,
puis réutilisée à chaque lancement.

\clearpage

\begin{figure}[H]
\centering
\resizebox{0.96\textwidth}{!}{%
\begin{tikzpicture}[
  font=\footnotesize,
  block/.style={
    draw=PrimaryBlue, fill=PrimaryBlue!8,
    rounded corners=7pt, align=center,
    minimum width=3.35cm,
    minimum height=0.95cm,
    font=\footnotesize\bfseries
  },
  module/.style={
    draw=PrimaryBlue!70, fill=white,
    rounded corners=4pt, align=center,
    minimum width=4.70cm,
    minimum height=0.78cm,
    font=\footnotesize
  },
  privacy/.style={
    draw=AccentOrange, fill=AccentOrange!10,
    rounded corners=4pt, align=center,
    minimum width=4.70cm,
    minimum height=0.78cm,
    font=\footnotesize\bfseries
  },
  side/.style={
    draw=AccentGreen, fill=AccentGreen!8,
    rounded corners=6pt, align=center,
    minimum width=3.65cm,
    minimum height=0.95cm,
    font=\footnotesize
  },
  arrow/.style={
    -{Stealth[length=4mm, width=2.5mm]},
    line width=1.15pt, PrimaryBlue},
  sidearrow/.style={
    -{Stealth[length=3mm, width=1.9mm]},
    thick},
  label/.style={
    font=\scriptsize, color=DarkGray,
    fill=white, inner sep=2pt}
]

\node[block] (cam) at (0,0)
  {Vidéo source\\\scriptsize fichier local};

\node[module]  (m1) at (5.4, 2.35)
  {M1 Acquisition};
\node[module]  (m2) at (5.4, 1.38)
  {M2 Détection};
\node[privacy] (m4) at (5.4, 0.41)
  {M4 Anonymisation};
\node[module]  (m3) at (5.4,-0.56)
  {M3 Tracking};
\node[module]  (m5) at (5.4,-1.53)
  {M5 Décision};
\node[module]  (m6) at (5.4,-2.50)
  {M6 Dashboard};

\begin{scope}[on background layer]
\node[
  draw=PrimaryBlue, fill=PrimaryBlue!5,
  rounded corners=10pt,
  fit=(m1)(m2)(m4)(m3)(m5)(m6),
  inner xsep=0.55cm, inner ysep=0.75cm,
  label={[font=\footnotesize\bfseries,
    color=PrimaryBlue]above:
    Raspberry Pi 5 — traitement embarqué}
] (rpi) {};
\end{scope}

\node[side] (model) at (11.2,1.38)
  {\textbf{Modèle ONNX}\\YOLO26n FP32};
\node[block] (browser) at (11.2,-2.50)
  {Navigateur\\\scriptsize Dashboard local};

\draw[arrow] (cam.east) --
  node[label, above] {frames}
  (rpi.west |- m1);
\draw[arrow] (m1) -- (m2);
\draw[arrow] (m2) -- (m4);
\draw[arrow] (m4) -- (m3);
\draw[arrow] (m3) -- (m5);
\draw[arrow] (m5) -- (m6);
\draw[arrow] (m6.east) --
  node[label, above] {HTTP/MJPEG}
  (browser.west);
\draw[sidearrow, AccentGreen] (model.west)
  -- node[label, above] {inférence}
  (m2.east);

\node[
  draw=AccentOrange, fill=AccentOrange!6,
  rounded corners=5pt, align=center,
  text width=9.8cm, font=\scriptsize,
  below=0.85cm of rpi
]
{Traitement entièrement local :
aucune image brute ne quitte
le Raspberry Pi 5.};

\end{tikzpicture}}
\caption{Architecture physique du système
         \textit{(Source : auteur)}}
\label{fig:physical_architecture}
\end{figure}

\subsection{Ordre du pipeline}

L'ordre M1→M2→M4→M3→M5→M6 n'est pas
arbitraire. La position de M4 avant M3
est une décision architecturale directement
liée à la Loi 09-08.

\begin{table}[H]
\centering
\caption{Description des modules du pipeline}
\footnotesize
\renewcommand{\arraystretch}{1.15}
\begin{tabularx}{\textwidth}{@{}
  >{\raggedright\arraybackslash}p{3.2cm}
  >{\raggedright\arraybackslash}p{2.4cm}
  >{\raggedright\arraybackslash}X
  >{\raggedright\arraybackslash}p{2.7cm}
@{}}
\toprule
\textbf{Module} &
\textbf{Entrée} &
\textbf{Traitement} &
\textbf{Sortie} \\
\midrule
M1 — Acquisition &
Source vidéo &
Lecture frame par frame via OpenCV &
Frame BGR \\
M2 — Détection &
Frame &
Inférence YOLO26n ONNX FP32 &
BBoxes + Classes \\
M4 — Anonymisation &
Frame + BBoxes &
Floutage gaussien classe \textit{person} &
Frame anonymisée \\
M3 — Tracking &
BBoxes &
IoU matching + comptage par zone &
Densités \\
M5 — Décision &
Densités + Types &
Scoring MDP $\pi^*(s)$ &
Phase + Durée \\
M6 — Dashboard &
Frame + Phase &
Streaming MJPEG + Flask &
Interface web \\
\bottomrule
\end{tabularx}
\label{tab:modules}
\end{table}

\subsection{Pourquoi M4 avant M3}

La question mérite d'être expliquée
clairement. On aurait pu appliquer
l'anonymisation après le tracking —
beaucoup de systèmes le font. On a
choisi de l'appliquer avant, pour
deux raisons.

La première est juridique. La Loi 09-08
impose de protéger les données personnelles
dès leur collecte, pas après traitement.
Même si aucun utilisateur ne voit les
frames intermédiaires, laisser circuler
des pixels identifiables entre modules
internes va à l'encontre de ce principe.

La deuxième est technique. Notre tracker
M3 est basé sur l'IoU — il travaille
uniquement avec les coordonnées des
boîtes englobantes $(x, y, w, h)$. Il
ne lit jamais les pixels à l'intérieur
des boîtes. Flouter la frame avant le
tracking ne dégrade donc pas les
performances de suivi d'un seul pixel.

\begin{figure}[H]
\centering
\resizebox{0.92\textwidth}{!}{%
\begin{tikzpicture}[
  box/.style={
    draw, rounded corners=4pt,
    align=center, font=\footnotesize,
    minimum width=3.3cm,
    minimum height=0.85cm},
  raw/.style={box,
    draw=red!65, fill=red!8},
  safe/.style={box,
    draw=AccentGreen,
    fill=AccentGreen!8},
  neutral/.style={box,
    draw=PrimaryBlue,
    fill=PrimaryBlue!7},
  note/.style={
    align=center,
    font=\footnotesize},
  arrow/.style={
    -{Stealth[length=4.5mm, width=3mm]},
    line width=1.4pt}
]

\node[note, font=\bfseries, color=red!70]
  at (0,0) {Ordre non retenu};
\node[neutral] (bm2) at (0,-1.0)
  {M2 Détection};
\node[raw] (bm3) at (0,-2.55)
  {M3 Tracking\\sur frame originale};
\node[raw] (bm4) at (0,-4.10)
  {M4 Anonymisation\\après tracking};
\node[note, color=red!70] at (0,-5.25)
  {La frame identifiable circule\\
  jusqu'au tracker};

\draw[arrow, red!65] (bm2) -- (bm3);
\draw[arrow, red!65] (bm3) -- (bm4);

\node[note, font=\bfseries,
  color=AccentGreen]
  at (8.4,0) {Ordre retenu};
\node[neutral] (gm2) at (8.4,-1.0)
  {M2 Détection};
\node[safe] (gm4) at (8.4,-2.55)
  {M4 Anonymisation\\immédiate};
\node[safe] (gm3) at (8.4,-4.10)
  {M3 Tracking IoU\\sur frame floutée};
\node[note, color=AccentGreen]
  at (8.4,-5.25)
  {Après M4, aucun pixel\\
  identifiable ne circule};

\draw[arrow, AccentGreen] (gm2) -- (gm4);
\draw[arrow, AccentGreen] (gm4) -- (gm3);

\node[
  draw=DarkGray, fill=LightGray,
  rounded corners=4pt,
  minimum width=2.6cm,
  minimum height=1.0cm,
  align=center, font=\footnotesize
] at (4.2,-6.4)
  {Loi 09-08\\Privacy-by-Design};

\end{tikzpicture}}
\caption{Justification de l'ordre M4 avant M3
         \textit{(Source : auteur)}}
\label{fig:order_justification}
\end{figure}

\section{Module M1 — Acquisition vidéo}

M1 est le point d'entrée du pipeline.
Il lit la vidéo source frame par frame
et alimente M2. L'implémentation repose
sur OpenCV, qui offre une API unifiée
pour la lecture de fichiers vidéo avec
une empreinte mémoire minimale sur
CPU ARM.

\begin{lstlisting}[language=Python,
  caption={Module M1 — lecture vidéo}]
cap = cv2.VideoCapture(source)
ret, frame = cap.read()
\end{lstlisting}

\section{Module M2 — Détection YOLO26n}

M2 reçoit chaque frame et produit les
boîtes englobantes avec leur classe
et leur score de confiance. Le benchmark
décrit au Chapitre~\ref{ch:implementation}
a conduit à retenir YOLO26n en format
ONNX FP32, pour les raisons détaillées
dans le Chapitre~\ref{ch:etat_art}.

Les tests INT8 ont montré une confusion
accrue entre les classes car, bus et
truck sur les frames de nuit et sur
les objets partiellement occultés.
Le format FP32 est donc conservé malgré
une empreinte plus grande.

\begin{equation}
\texttt{yolo26n.pt}
\xrightarrow{\text{export ONNX}}
\texttt{YOLO26n\_step3\_960\_best.onnx}
\label{eq:export}
\end{equation}

\section{Module M4 — Anonymisation
         Privacy-by-Design}
\label{sec:m4}

\subsection{Positionnement et rôle}

M4 est le module le plus sensible du
point de vue juridique. Il reçoit la
frame originale et les boîtes de M2,
applique le floutage sur toutes les
boîtes de classe \textit{person}, et
transmet une frame anonymisée à M3.
Aucun pixel identifiable ne circule
au-delà de ce module.

\subsection{Choix du noyau gaussien}

La valeur $51 \times 51$ pixels a été
choisie par ajustement visuel sur les
frames du dashboard. L'objectif était
d'obtenir un floutage visiblement
protecteur sans rendre la zone
complètement illisible.

Un noyau $25 \times 25$ donnait un
résultat parfois insuffisant sur les
visages proches de la caméra. Un noyau
$75 \times 75$ était efficacement
protecteur mais dégradait la lisibilité
des zones — notamment quand plusieurs
piétons se chevauchent. La valeur
$51 \times 51$ représente le bon
compromis, validé visuellement sur
plusieurs frames de référence.

\begin{algorithm}[H]
\caption{Anonymisation M4 —
         Privacy-by-Design}
\begin{algorithmic}[1]
\REQUIRE frame $F$, détections $D$
\ENSURE frame anonymisée $F'$
\STATE $F' \leftarrow F$
\FOR{chaque détection $d \in D$}
  \IF{$d.\text{class} = \text{person}$}
    \STATE $x_1, y_1, x_2, y_2
      \leftarrow d.\text{bbox}$
    \STATE $\text{roi} \leftarrow
      F'[y_1:y_2,\ x_1:x_2]$
    \STATE $\text{roi} \leftarrow
      \text{GaussianBlur}(
      \text{roi}, (51, 51), 0)$
    \STATE $F'[y_1:y_2,\ x_1:x_2]
      \leftarrow \text{roi}$
  \ENDIF
\ENDFOR
\RETURN $F'$
\end{algorithmic}
\end{algorithm}

\subsection{Compatibilité avec M3}

Le tracker IoU travaille uniquement
avec les coordonnées des boîtes
$(x, y, w, h)$. Il ne lit jamais
les pixels à l'intérieur. Flouter
la frame avant le tracking n'affecte
donc en rien les performances de suivi.

\begin{table}[H]
\centering
\caption{Compatibilité de M4 avec
         les modules suivants}
\begin{tabularx}{\textwidth}{lXl}
\toprule
\textbf{Module} &
\textbf{Données utilisées} &
\textbf{Impacté par M4 ?} \\
\midrule
M3 Tracker IoU &
Coordonnées bbox uniquement &
Non \\
M5 MDP &
Comptages par zone &
Non \\
M6 Dashboard &
Frame anonymisée affichée &
Non \\
\bottomrule
\end{tabularx}
\label{tab:m4_compat}
\end{table}

\subsection{Pourquoi pas ByteTrack}
\label{sec:bytetrack}

ByteTrack offre une meilleure robustesse
aux occlusions grâce à un modèle de
ré-identification visuelle qui extrait
des caractéristiques du corps pour
associer les détections entre frames.
Ce mécanisme est incompatible avec
notre architecture pour deux raisons.

Juridiquement, extraire des
caractéristiques visuelles d'une personne
constitue un traitement de données
biométriques au sens de la Loi 09-08,
soumis à des exigences de consentement
que notre système de surveillance ne
peut pas satisfaire.

Matériellement, le modèle de
ré-identification additionnel de ByteTrack
consomme des ressources CPU incompatibles
avec le Raspberry Pi 5 en mode temps réel.

Le tracker IoU retenu règle les deux
problèmes : il n'utilise que les
coordonnées, sans jamais lire les pixels.

\section{Module M3 — Tracking IoU
         et comptage par zones}

\subsection{Rôle}

M3 reçoit les frames anonymisées de M4
et les boîtes de M2. Il fait deux choses :
associer chaque détection à un identifiant
de track persistant d'une frame à l'autre,
et affecter chaque objet à une zone
directionnelle pour produire les densités
transmises à M5.

\subsection{Algorithme IoU}

Pour chaque détection à l'instant $t$,
le tracker cherche la boîte de l'instant
$t-1$ qui maximise le chevauchement :

\begin{equation}
\text{track}^* =
\arg\max_{j}\ \text{IoU}(
\hat{B}_i^t,\ \hat{B}_j^{t-1})
\label{eq:iou_tracking}
\end{equation}

Une association est validée si
$\text{IoU} \geq 0.3$. En dessous,
un nouveau track est créé. Un track
sans association depuis $N_{lost} = 5$
frames consécutives est supprimé.

\begin{algorithm}[H]
\caption{Tracking IoU — M3}
\begin{algorithmic}[1]
\REQUIRE détections $D^t$,
         tracks actifs $T^{t-1}$
\ENSURE tracks $T^t$, densités
\STATE $M_{ij} \leftarrow
  \text{IoU}(D_i^t, T_j^{t-1})$
\STATE Résoudre l'assignation optimale
  (algorithme hongrois)
\FOR{paire associée $(i,j)$,
     $M_{ij} \geq 0.3$}
  \STATE Mettre à jour track $j$
\ENDFOR
\FOR{détection $i$ non associée}
  \STATE Créer nouveau track
\ENDFOR
\FOR{track $j$ non associé}
  \STATE $\text{lost}_j \mathrel{+}= 1$
  \IF{$\text{lost}_j > 5$}
    \STATE Supprimer track $j$
  \ENDIF
\ENDFOR
\RETURN $T^t$
\end{algorithmic}
\end{algorithm}

\subsection{Zones polygonales}

Les zones directionnelles sont des
polygones définis manuellement sur
une frame de référence. La calibration
se fait à la souris via l'outil
\texttt{calibrate\_interactive.py} :
on ouvre une frame de la vidéo, on
clique les coins de chaque zone, et
les coordonnées sont sauvegardées dans
un fichier JSON réutilisé à chaque
lancement du pipeline. Cette calibration
manuelle est nécessaire parce que la
géométrie d'une intersection varie
selon l'angle et la hauteur de la
caméra — il n'existe pas de configuration
automatique universelle.

\begin{table}[H]
\centering
\caption{Zones polygonales de détection}
\begin{tabularx}{\textwidth}{llXl}
\toprule
\textbf{Zone} &
\textbf{Variable} &
\textbf{Description} &
\textbf{Couleur} \\
\midrule
Nord    & $q_N$ &
Véhicules arrivant du nord & Bleu \\
Sud     & $q_S$ &
Véhicules arrivant du sud & Vert \\
Est     & $q_E$ &
Véhicules arrivant de l'est & Jaune \\
Ouest   & $q_O$ &
Véhicules arrivant de l'ouest & Orange \\
Piétons & $p$   &
Zone centrale de passage & Magenta \\
\bottomrule
\end{tabularx}
\label{tab:zones}
\end{table}

L'affectation d'un objet à une zone
utilise le centre bas de sa boîte
$(x_c, y_{bas})$, qui correspond
approximativement au point de contact
au sol. C'est plus robuste que le centre
géométrique, qui se déplace quand
l'objet est partiellement occulté.

\begin{equation}
\text{zone}(B) = z \in Z\
\text{tel que}\
(x_c, y_{bas}) \in \text{polygone}(z)
\label{eq:zone_assign}
\end{equation}

\subsection{Lissage par fenêtre glissante}

Les densités transmises à M5 sont
calculées sur une fenêtre glissante
de $W = 10$ frames :

\begin{equation}
q_z = \frac{1}{W}
\sum_{t=T-W}^{T} \text{count}(z, t)
\label{eq:density}
\end{equation}

La valeur $W = 10$ est un compromis
empirique entre réactivité et stabilité.
Une fenêtre de 1 ou 3 frames rend le
système trop sensible aux faux positifs
transitoires — une seule détection
erronée peut changer la phase de feux.
Une fenêtre de 30 frames stabilise les
scores mais ralentit la réponse aux
vrais changements du trafic. La valeur
10 frames donne une décision qui réagit
aux tendances plutôt qu'aux fluctuations.

\subsection{Comptage par classe et par zone}

Le tracker distingue les classes dans
chaque zone pour alimenter M5 avec
un état détaillé :

\begin{table}[H]
\centering
\caption{Exemple d'état par classe et par zone}
\begin{tabularx}{\textwidth}{lcccccc}
\toprule
\textbf{Zone} &
\textbf{car} & \textbf{bus} &
\textbf{truck} & \textbf{moto} &
\textbf{emergency} & \textbf{person} \\
\midrule
Nord    & 8 & 1 & 0 & 2 & 0 & 0 \\
Sud     & 2 & 0 & 1 & 1 & 0 & 0 \\
Est     & 4 & 0 & 0 & 3 & 1 & 0 \\
Ouest   & 3 & 0 & 0 & 1 & 0 & 0 \\
Piétons & 0 & 0 & 0 & 0 & 0 & 3 \\
\midrule
\textbf{Total} &
17 & 1 & 1 & 7 & 1 & 3 \\
\bottomrule
\end{tabularx}
\label{tab:zone_counts}
\end{table}

Dans cet exemple, le vecteur d'état
transmis à M5 est $s = (11, 4, 8, 4, 3, 1)$.
Comme $e = 1$, la phase Emergency est
activée immédiatement, quelle que soit
la valeur des autres scores.

\section{Module M5 — Décision MDP}

M5 reçoit le vecteur d'état de M3 et
calcule la phase de feux à activer.
La politique est déterministe — pas
de réseau neuronal, pas d'apprentissage
en ligne. La décision se réduit à
quelques comparaisons et un calcul
de score, exécutable en microsecondes
sur le CPU ARM.

\subsection{Formulation}

L'état courant :

\begin{equation}
s = (q_N,\ q_S,\ q_E,\ q_O,\ p,\ e)
\label{eq:state_m5}
\end{equation}

Le score pour chaque direction $i$ :

\begin{equation}
\text{Score}(i) =
w_v \cdot q_i^{veh}
+ w_p \cdot q_i^{ped}
+ w_e \cdot q_i^{urg}
\label{eq:scoring}
\end{equation}

\begin{table}[H]
\centering
\caption{Coefficients et seuils du module M5}
\begin{tabularx}{\textwidth}{llXl}
\toprule
\textbf{Param.} &
\textbf{Valeur} &
\textbf{Logique} &
\textbf{Classes} \\
\midrule
$w_v$ & 1.0 &
Référence de base &
car, bus, truck, moto \\
$w_p$ & 1.5 &
Piétons prioritaires, pas
au point de bloquer le trafic &
person \\
$w_e$ & 100.0 &
Valeur assez haute pour
qu'aucun autre score ne
puisse la dépasser &
emergency\_vehicle \\
$\theta_H$ & 10 veh &
Seuil densité haute → 45s &
— \\
$\theta_M$ & 5 veh &
Seuil densité moyenne → 30s &
— \\
\bottomrule
\end{tabularx}
\label{tab:mdp_params}
\end{table}

La politique :

\begin{equation}
\pi^*(s) =
\begin{cases}
\text{Emergency} &
\text{si } e > 0 \\[5pt]
\text{Pedestrian} &
\text{si } p > 0 \wedge
\max_i q_i^{veh} < \theta_H \\[5pt]
\arg\max_i \text{Score}(i) &
\text{sinon}
\end{cases}
\label{eq:policy_m5}
\end{equation}

La durée de phase :

\begin{equation}
D^* =
\begin{cases}
45\text{s} &
q_i^{veh} \geq \theta_H \\
30\text{s} &
\theta_M \leq q_i^{veh} < \theta_H \\
15\text{s} &
q_i^{veh} < \theta_M
\end{cases}
\label{eq:duration_m5}
\end{equation}


\subsection{Exemple numérique — Frame F181}

Pour illustrer concrètement le fonctionnement
de la politique $\pi^*(s)$, prenons l'état
réel capturé à la frame F181 de la vidéo
Bellevue NE8th :

\begin{equation}
s = (q_N=12,\ q_S=0,\ q_E=2,\ q_W=6,\
p=7,\ e=0)
\label{eq:exemple_etat}
\end{equation}

\textbf{Étape 1 — Vérification urgence :}
$e = 0$, pas de véhicule d'urgence.
On passe à l'étape suivante.

\textbf{Étape 2 — Vérification piétons :}
$p = 7 > 0$ mais $\max(q_i^{veh}) =
\max(12, 0, 2, 6) = 12 \geq \theta_H = 10$.
Le trafic est dense — la phase Pedestrian
n'est pas déclenchée.

\textbf{Étape 3 — Calcul des scores :}

\begin{align}
\text{Score}(NS) &= w_v \cdot (q_N + q_S)
+ w_p \cdot p \notag \\
&= 1.0 \times (12 + 0)
+ 1.5 \times 7 \notag \\
&= 12.0 + 10.5 = \mathbf{22.5}
\label{eq:score_ns_exemple}
\end{align}

\begin{align}
\text{Score}(EW) &= w_v \cdot (q_E + q_W)
+ w_p \cdot p \notag \\
&= 1.0 \times (2 + 6)
+ 1.5 \times 7 \notag \\
&= 8.0 + 10.5 = \mathbf{18.5}
\label{eq:score_ew_exemple}
\end{align}

\textbf{Étape 4 — Décision :}
$\text{Score}(NS) = 22.5 >
\text{Score}(EW) = 18.5$

\begin{equation}
\pi^*(s) = \text{North-South}
\label{eq:decision_exemple}
\end{equation}

\textbf{Étape 5 — Durée :}
$q_N = 12 \geq \theta_H = 10$

\begin{equation}
D^* = 45\text{s}
\label{eq:duree_exemple}
\end{equation}

\begin{table}[H]
\centering
\caption{Récapitulatif — Exemple numérique
         Frame F181}
\begin{tabularx}{\textwidth}{lX}
\toprule
\textbf{Élément} & \textbf{Valeur} \\
\midrule
État détecté &
$q_N=12, q_S=0, q_E=2, q_W=6,
p=7, e=0$ \\
Score NS & 22.5 \\
Score EW & 18.5 \\
Phase sélectionnée & North-South \\
Durée & 45s \\
Raison & Axe Nord/Sud dominant,
densité Nord élevée ($q_N=12
\geq \theta_H=10$) \\
\bottomrule
\end{tabularx}
\label{tab:exemple_numerique_mdp}
\end{table}

Cet exemple illustre trois propriétés
importantes de la politique. Premièrement,
les piétons contribuent aux deux scores
via $w_p = 1.5$ — ils n'ont pas de phase
dédiée ici car le trafic est dense.
Deuxièmement, même avec $q_S = 0$, l'axe
NS est sélectionné car $q_N = 12$ domine
largement. Troisièmement, la durée maximale
de 45s est attribuée car la densité dépasse
le seuil $\theta_H = 10$.


\section{Module M6 — Dashboard Flask}

\subsection{Contexte de développement}

Le dashboard n'a pas été conçu d'un coup.
Il a évolué par itérations. La première
version affichait un flux MJPEG simple,
les statistiques JSON et les scores MDP.
Ensuite ont été ajoutés : la vidéo source
en parallèle de l'image analysée, le
bouton d'analyse fixe, l'affichage des
polygones de zones, la suppression des
textes inutiles superposés sur l'image,
et l'amélioration progressive de l'interface.
La version finale sépare clairement
la preuve visuelle IA et les données
décisionnelles dans quatre panneaux.

\subsection{Architecture technique}

Le dashboard est implémenté avec Flask
(Python), accessible via navigateur
sur le réseau local. L'architecture est
thread-safe : le pipeline M1→M5 tourne
dans un thread dédié pendant que Flask
répond aux requêtes HTTP dans un thread
séparé. L'image annotée est transmise
en MJPEG sur le port 5000 :

\begin{center}
\texttt{http://192.168.1.117:5000}
\end{center}

\begin{table}[H]
\centering
\caption{Endpoints du dashboard Flask}
\begin{tabularx}{\textwidth}{llX}
\toprule
\textbf{Route} &
\textbf{Méthode} &
\textbf{Description} \\
\midrule
\texttt{/} &
GET & Page principale \\
\texttt{/api/state} &
GET & État JSON temps réel \\
\texttt{/api/analyze} &
POST & Analyser la frame courante \\
\texttt{/api/seek} &
POST & Aller à une seconde de la vidéo \\
\texttt{/source\_video} &
GET & Flux vidéo source \\
\texttt{/analysis\_frame.jpg} &
GET & Image analysée annotée \\
\texttt{/health} &
GET & Vérification du serveur \\
\bottomrule
\end{tabularx}
\label{tab:flask_endpoints}
\end{table}

\subsection{Interface — quatre panneaux}

L'interface est organisée en quatre zones.

Le \textbf{panneau image IA} affiche la
frame anonymisée avec les boîtes de
détection colorées par classe, les
identifiants de tracks, les polygones
de zones en transparence et les zones
floutées sur les piétons. Aucune
statistique n'est superposée sur l'image
pour conserver sa lisibilité.

Le \textbf{panneau KPI} affiche en temps
réel les métriques opérationnelles du
pipeline : FPS, latence d'analyse,
nombre d'objets détectés, piétons
anonymisés, tracks actifs et numéro
de frame.

Le \textbf{panneau décision MDP} présente
la phase active, sa durée, les scores
NS/EO et l'état des feux par direction
avec des indicateurs colorés rouge/vert.

Le \textbf{panneau zones} détaille la
composition du trafic par zone directionnelle
— nombre de véhicules par classe, piétons,
score MDP et indicateur de feu.

\subsection{Conformité Privacy-by-Design}

Le dashboard respecte les contraintes
de la Loi 09-08 à quatre niveaux : seule
la frame anonymisée est transmise et
affichée, aucune frame n'est sauvegardée
sur disque, l'accès est limité au réseau
local sans exposition internet, et l'API
\texttt{/api/state} ne retourne que des
comptages agrégés.

\section{Synthèse des choix technologiques}

\begin{table}[H]
\centering
\caption{Récapitulatif des choix par module}
\begin{tabularx}{\textwidth}{llXl}
\toprule
\textbf{Module} &
\textbf{Technologie} &
\textbf{Raison principale} &
\textbf{Réf.} \\
\midrule
M1 & OpenCV &
API unifiée, léger ARM &
\cite{ficili2025} \\
M2 & YOLO26n ONNX FP32 &
FPS maximal, NMS-Free &
\cite{sapkota2025} \\
M4 & GaussianBlur 51×51 &
Visiblement protecteur,
compatible M3 &
\cite{jeremiah2024} \\
M3 & IoU Tracker &
Léger, sans features visuelles,
conforme Loi 09-08 &
\cite{wang2025} \\
M5 & MDP scoring &
Déterministe, temps réel &
\cite{michailidis2025} \\
M6 & Flask + MJPEG &
Python natif, streaming léger &
\cite{ficili2025} \\
\bottomrule
\end{tabularx}
\label{tab:tech_choices}
\end{table}

\section*{Conclusion}
\addcontentsline{toc}{section}{Conclusion}

Ce chapitre a décrit l'architecture du
système et justifié chaque choix de
conception. La décision centrale —
placer M4 avant M3 — découle directement
de la nature position-based du tracker
IoU et de l'exigence Privacy-by-Design.
Les autres choix, du noyau gaussien
$51 \times 51$ à la fenêtre glissante
de 10 frames, ont tous une justification
empirique documentée.

Le chapitre suivant détaille la construction
du dataset et le benchmark expérimental
qui ont validé ces choix.